% GENERATED FILE. Edit canonical lessons, then run npm run generate:books.
% canonical-source-hash: fnv1a64:2a5c97a453a20094
% canonical-lessons: LA-C463-holitor, LA-C463-piscarius, LA-C463-sporta, LA-C463-corbis, LA-C463-canistrum

\chapter{Vegetable grower, Fishmonger, Plaited basket, Wicker basket, Bread basket}
\label{ch:la-vegetable-grower-fishmonger-plaited-basket-wicker-basket-bread-basket}
\hlchaptermodality{463}

\begin{chapteropening}
\textbf{By the end of this chapter:} \emph{I can talk about market sellers and what they carry things in: the vegetable grower, the fishmonger, and three kinds of basket.}

Complete the last lesson of chapter 463: I can talk about market sellers and what they carry things in: the vegetable grower, the fishmonger, and three kinds of basket.
\end{chapteropening}

\section[holitor, holitōris]{holitor, holitōris — a vegetable grower, a market gardener}

\label{lesson:LA-C463-holitor}



\begin{quote}
Before the new one: say the Latin for goods, then the Latin for a hawker.
\end{quote}

\subsection*{You'll want to know: holitor, holitōris}
\textbf{holitor, holitōris} — \textquotedblleft{}a vegetable grower, a market gardener\textquotedblright{}.

\begin{quote}
Holitor holera in forum portat et vendit.
\end{quote}

\begin{quote}
\emph{The market gardener carries vegetables to the forum and sells them.}
\end{quote}

Third declension, masculine: \emph{holitor, holitōris m.}, also spelled \emph{olitor}. From \emph{holus} 'vegetables'.

\begin{cousinweb}
\emph{holitor} is built on \emph{holus, holeris} 'vegetables, greens' with the suffix \emph{-tor} 'one who' (Lewis \& Short, under \emph{olitor}).
\end{cousinweb}

\begin{grammarlens}[title={What you've built}]
One more word for telling, explaining and describing at B1.
\end{grammarlens}

\subsection*{Your turn}
\begin{itemize}
\raggedright
  \item \emph{Say it:} \emph{holitor, holitōris}
  \item \emph{Say it:} \emph{holitor, holitōris}, once more
  \item \emph{Say it:} a sentence of your own with \emph{holitor, holitōris}
  \item \emph{Recall:} say the Latin for goods, then the Latin for a hawker, then say \emph{holitor, holitōris} again
\end{itemize}

\begin{tcolorbox}[breakable,colback=teal!4,colframe=teal!35!black,title={Before you move on}]
What does holitor, holitōris mean? (\textquotedblleft{}A vegetable grower, a market gardener\textquotedblright{}.) Say it once more.
\end{tcolorbox}
\section[piscārius, piscāriī]{piscārius, piscāriī — a fishmonger}

\label{lesson:LA-C463-piscarius}



\begin{quote}
Before the new one: say the Latin for a hawker, then the Latin for a vegetable grower.
\end{quote}

\subsection*{You'll want to know: piscārius, piscāriī}
\textbf{piscārius, piscāriī} — \textquotedblleft{}a fishmonger\textquotedblright{}.

\begin{quote}
Piscārius piscēs recentēs clāmat.
\end{quote}

\begin{quote}
\emph{The fishmonger calls out his fresh fish.}
\end{quote}

Second declension, masculine: \emph{piscārius, piscāriī m.} Also an adjective: \emph{forum piscārium}, the fish market.

\begin{cousinweb}
\emph{piscārius} is built on \emph{piscis} 'fish' with the suffix \emph{-ārius} (EtymDB; Lewis \& Short); \emph{piscis} comes from PIE \emph{peys\'{k}-}, and English \emph{fish} is a Germanic cousin.
\end{cousinweb}

\begin{grammarlens}[title={What you've built}]
One more word for telling, explaining and describing at B1.
\end{grammarlens}

\subsection*{Your turn}
\begin{itemize}
\raggedright
  \item \emph{Say it:} \emph{piscārius, piscāriī}
  \item \emph{Say it:} \emph{piscārius, piscāriī}, once more
  \item \emph{Say it:} a sentence of your own with \emph{piscārius, piscāriī}
  \item \emph{Recall:} say the Latin for a hawker, then the Latin for a vegetable grower, then say \emph{piscārius, piscāriī} again
\end{itemize}

\begin{tcolorbox}[breakable,colback=teal!4,colframe=teal!35!black,title={Before you move on}]
What does piscārius, piscāriī mean? (\textquotedblleft{}A fishmonger\textquotedblright{}.) Say it once more.
\end{tcolorbox}
\section[sporta, sportae]{sporta, sportae — a plaited basket, a hamper}

\label{lesson:LA-C463-sporta}



\begin{quote}
Before the new one: say the Latin for a vegetable grower, then the Latin for a fishmonger.
\end{quote}

\subsection*{You'll want to know: sporta, sportae}
\textbf{sporta, sportae} — \textquotedblleft{}a plaited basket, a hamper\textquotedblright{}.

\begin{quote}
Ancilla pānem et cāseum in sportā portat.
\end{quote}

\begin{quote}
\emph{The maid carries bread and cheese in a hamper.}
\end{quote}

First declension, feminine: \emph{sporta, sportae f.} Its diminutive \emph{sportula} was the food basket, later the cash, a patron gave his clients.

\begin{cousinweb}
\emph{sporta} is usually traced to Greek \emph{spyris} 'plaited basket' (Lewis \& Short compare it); EtymDB links it with PIE \emph{sper-} 'to twist, to turn'. The path is not settled.
\end{cousinweb}

\begin{grammarlens}[title={What you've built}]
One more word for telling, explaining and describing at B1.
\end{grammarlens}

\subsection*{Your turn}
\begin{itemize}
\raggedright
  \item \emph{Say it:} \emph{sporta, sportae}
  \item \emph{Say it:} \emph{sporta, sportae}, once more
  \item \emph{Say it:} a sentence of your own with \emph{sporta, sportae}
  \item \emph{Recall:} say the Latin for a vegetable grower, then the Latin for a fishmonger, then say \emph{sporta, sportae} again
\end{itemize}

\begin{tcolorbox}[breakable,colback=teal!4,colframe=teal!35!black,title={Before you move on}]
What does sporta, sportae mean? (\textquotedblleft{}A plaited basket, a hamper\textquotedblright{}.) Say it once more.
\end{tcolorbox}
\section[corbis, corbis]{corbis, corbis — a wicker basket}

\label{lesson:LA-C463-corbis}



\begin{quote}
Before the new one: say the Latin for a fishmonger, then the Latin for a plaited basket.
\end{quote}

\subsection*{You'll want to know: corbis, corbis}
\textbf{corbis, corbis} — \textquotedblleft{}a wicker basket\textquotedblright{}.

\begin{quote}
Servī ūvās in corbēs pōnunt.
\end{quote}

\begin{quote}
\emph{The slaves put the grapes into wicker baskets.}
\end{quote}

Third declension, masculine or feminine: \emph{corbis, corbis}. A big basket of woven twigs, for the field or the vineyard.

\begin{cousinweb}
The origin of \emph{corbis} is not settled: an ancient writer explained it from the bent (\emph{curvātae}) twigs it is woven of, while Lewis \& Short compare Greek \emph{kolpos} 'fold'. EtymDB gives no parent.
\end{cousinweb}

\begin{grammarlens}[title={What you've built}]
One more word for telling, explaining and describing at B1.
\end{grammarlens}

\subsection*{Your turn}
\begin{itemize}
\raggedright
  \item \emph{Say it:} \emph{corbis, corbis}
  \item \emph{Say it:} \emph{corbis, corbis}, once more
  \item \emph{Say it:} a sentence of your own with \emph{corbis, corbis}
  \item \emph{Recall:} say the Latin for a fishmonger, then the Latin for a plaited basket, then say \emph{corbis, corbis} again
\end{itemize}

\begin{tcolorbox}[breakable,colback=teal!4,colframe=teal!35!black,title={Before you move on}]
What does corbis, corbis mean? (\textquotedblleft{}A wicker basket\textquotedblright{}.) Say it once more.
\end{tcolorbox}
\section[canistrum, canistrī]{canistrum, canistrī — a bread basket, a reed basket}

\label{lesson:LA-C463-canistrum}



\begin{quote}
Before the new one: say the Latin for a plaited basket, then the Latin for a wicker basket.
\end{quote}

\subsection*{You'll want to know: canistrum, canistrī}
\textbf{canistrum, canistrī} — \textquotedblleft{}a bread basket, a reed basket\textquotedblright{}.

\begin{quote}
In cēnā canistrum plēnum pānis est.
\end{quote}

\begin{quote}
\emph{At dinner there is a basket full of bread.}
\end{quote}

Second declension, neuter: \emph{canistrum, canistrī n.}, mostly plural. English \emph{canister} comes from it.

\begin{cousinweb}
\emph{canistrum} was borrowed from Greek \emph{kanastron} 'basket of reeds' (EtymDB; Lewis \& Short), from \emph{kanna} 'reed'. English \emph{canister} comes from it.
\end{cousinweb}

\begin{grammarlens}[title={What you've built}]
One more word for telling, explaining and describing at B1.
\end{grammarlens}

\subsection*{Your turn}
\begin{itemize}
\raggedright
  \item \emph{Say it:} \emph{canistrum, canistrī}
  \item \emph{Say it:} \emph{canistrum, canistrī}, once more
  \item \emph{Say it:} a sentence of your own with \emph{canistrum, canistrī}
  \item \emph{Recall:} say the Latin for a plaited basket, then the Latin for a wicker basket, then say \emph{canistrum, canistrī} again
\end{itemize}

\begin{tcolorbox}[breakable,colback=teal!4,colframe=teal!35!black,title={Before you move on}]
What does canistrum, canistrī mean? (\textquotedblleft{}A bread basket, a reed basket\textquotedblright{}.) Say it once more.
\end{tcolorbox}
